% Section III — System
% Draft 2026-10-20. Code: backend/rag/*, backend/kg/queries.py, eval/vector_baseline.py.

\section{System}
\label{sec:system}

Fig.~\ref{fig:pipeline} shows the pipeline. Both systems compared in this paper
share every step except the choice of the papers given to the LLM.

\begin{figure}[t]
\centering
% Fig. 1 — pipeline. Needs \usepackage{tikz} and
% \usetikzlibrary{positioning,arrows.meta} in main.tex.
% Fits one IEEE column (about 8.8 cm).

\begin{tikzpicture}[
    font=\footnotesize,
    node distance=4mm,
    >={Stealth[length=1.6mm]},
    step/.style={draw, rounded corners=1.5pt, align=center,
                 inner sep=2.5pt, minimum height=6mm},
    wide/.style={step, text width=7.6cm},
    half/.style={step, text width=3.55cm, minimum height=3.1cm, fill=gray!15},
    io/.style={align=center, inner sep=1pt},
]
  % shared input
  \node[io] (q) {Question};
  \node[wide, below=3mm of q] (an)
    {\textbf{Question analysis:} type (rules, then LLM),\\
     entities (LLM), normalisation to ORKG forms};

  % the only step that differs between the two systems
  \node[half, anchor=north east] (gr) at ([xshift=-1.5mm, yshift=-6mm]an.south)
    {\textbf{Graph retrieval}\\[1pt]
     SPARQL strategies per type\\
     (entities, comparison, claim\\
     evidence, benchmark, origin,\\
     title) $\to$ role-aware\\
     ranking and filtering $\to$\\
     enrichment of the top 15};
  \node[half, anchor=north west] (de) at ([xshift=1.5mm, yshift=-6mm]an.south)
    {\textbf{Dense retrieval}\\[1pt]
     embed the question\\
     (\texttt{mxbai-embed-large})\\
     $\to$ cosine similarity\\
     with cards of all\\
     68,806 ORKG papers};

  % shared output
  \node[wide, anchor=north] (top) at ([yshift=-5mm]an.south |- gr.south)
    {\textbf{Top 8 papers}};
  \node[wide, below=of top] (ctx)
    {\textbf{Context:} full ORKG records fetched live,\\
     up to 15 \emph{predicate: value} rows per paper};
  \node[wide, below=of ctx] (llm)
    {\textbf{LLM} with a prompt per question type\\
     (local via Ollama, or a cloud API)};
  \node[io, below=3mm of llm] (ans)
    {Answer citing papers by exact title\\
     (claims: SUPPORTED / REFUTED / INSUFFICIENT)};

  % arrows
  \draw[->] (q) -- (an);
  \draw[->] (an.south -| gr.north) -- (gr.north);
  \draw[->, dashed] (an.south -| de.north) -- (de.north)
    node[midway, right, font=\scriptsize] {question text};
  \draw[->] (gr.south) -- (gr.south |- top.north);
  \draw[->] (de.south) -- (de.south |- top.north);
  \draw[->] (top) -- (ctx);
  \draw[->] (ctx) -- (llm);
  \draw[->] (llm) -- (ans);
\end{tikzpicture}

\caption{The question-answering pipeline. Only the retrieval step (shaded)
differs between the two systems; dense retrieval uses only the question text
(dashed).}
\label{fig:pipeline}
\end{figure}

\subsection{Question Analysis}

The question is first assigned one of six types (Section~\ref{sec:method}).
Unambiguous forms are decided by rules (e.g.\ ``find the paper titled~\ldots'',
``compare~\ldots{} versus~\ldots'', a yes/no question stating a claim); the rest
are classified by the LLM. In parallel, the LLM extracts typed entities:
methods, datasets, tasks, research fields and metrics. If none are found,
keyword phrases from the question are used instead. A normalisation step
expands each entity to the surface forms used in ORKG, e.g.\ \emph{CNN} to
\emph{convolutional neural network}, and singular and plural forms.

\subsection{Graph Retrieval}

Graph retrieval queries the public ORKG SPARQL endpoint and uses no embeddings.
The question type selects a set of query strategies, each following a
different path in the graph from a paper through its contributions:

\begin{itemize}
\item \textbf{Entity paths}: contributions whose method, dataset or research
  problem matches an entity, and their combination when both a method and a
  dataset are given;
\item \textbf{Comparison}: papers whose contributions contain both compared
  methods;
\item \textbf{Claim evidence}: contributions covering most terms of the
  question, returned with the predicate that links each value (e.g.\
  \emph{baseline} versus \emph{proposed model});
\item \textbf{Benchmark traversal}: from a contribution through its model and
  benchmark nodes to the evaluated dataset and the reported scores;
\item \textbf{Origin lookup}: the earliest papers recording a method, as
  candidates for the paper that introduced it;
\item \textbf{Title lookup}: exact and keyword title search for paper-lookup
  questions.
\end{itemize}

All planned queries run in parallel. Their results are reduced to one
profile per paper and scored by the roles in which the question's entities
occur. For example, a method in a method or baseline role scores higher than
the same string elsewhere, an exact title match dominates, and papers that
carry benchmark results get a bonus. Short terms must match whole words, so
\emph{GAN} does not match \emph{organ}. When a question names both a method and
a dataset, papers matching only one of them are filtered out. The 15
best-scored candidates are then \emph{enriched}: their full contribution
records are fetched in one query and the papers are re-scored on the complete
record. This step matters because the matching query returns only the matched
value, so a paper that compares ResNet with VGG may first arrive with its VGG
row only. The top 8 papers are passed on.

\subsection{Dense Retrieval Baseline}

The dense baseline replaces only the step above. Every paper in the ORKG RDF dump
(68,806 papers) is represented by one text card: its title, research field
and the labels of its contribution values, research problems first, at most
1,500 characters. The cards are embedded with \texttt{mxbai-embed-large};
the question is embedded with the model's retrieval instruction, and the 8
papers with the highest cosine similarity are returned. The cards use the
same graph data that graph retrieval can reach, so the comparison is between
two ways of selecting papers from the same information, not between different
corpora.

\subsection{Context and Generation}

For the 8 selected papers, from either retriever, the full ORKG records are
fetched live and rendered as one block per paper: the title followed by
\emph{predicate: value} rows, up to 15 rows per paper. The rows are chosen in order of priority: values matching the question's
entities, benchmark results, then values under key predicates (method,
baseline, dataset, metric, research problem). Each row keeps the predicate,
so the LLM sees the role of a value, e.g.\ \emph{Baseline: Faster R-CNN}. Claim questions get an evidence-oriented layout of the same
rows. Papers are identified by title rather than by number, which prevents
citations of non-existent sources. A prompt template per question type asks
the LLM to answer only from the context, to name papers by their exact title,
and, for claims, to give a verdict (SUPPORTED, REFUTED or INSUFFICIENT) with
evidence. The system runs with local models through Ollama or with cloud APIs,
behind one client interface.
